Ambiguity note: “additive basis of order $k$” is used in the literature both for an exact basis (every $n\in\mathbb N$ is a sum of $k$ elements) and for an asymptotic basis (all sufficiently large $n$ are such sums). The problem statement uses the common Erd\H{o}s-problems convention “basis of order $k$” in the asymptotic sense; I adopt that as the minimal corrected statement.

Definitions and conventions used below:
- $\mathbb N:=\{1,2,3,\dots\}$.
- For a set $A\subset\mathbb N$ and integer $k\ge 1$, write
  \[kA:=\{a_1+\cdots+a_k: a_i\in A\}.\]
  (Summands may repeat.)
- $A$ is an (asymptotic) additive basis of order $k$ if $\exists N_0\in\mathbb N$ such that $[N_0,\infty)\cap\mathbb N\subset kA$.
- $A$ is “minimal in the strong sense” if for every infinite $B\subset A$, the set $A\setminus B$ is not an additive basis of order $k$, i.e.
  \[\forall\, B\subset A\ (B\text{ infinite})\ \forall N\ \exists n\ge N\text{ with }n\notin k(A\setminus B).\]

Question (corrected statement): Given $k\ge 1$ and $A\subset\mathbb N$ an additive basis of order $k$ that is minimal in the above strong sense, must there exist an infinite $B\subset A$ such that $A\setminus B$ is an additive basis of order $k+1$?
